% hmm_mcq_challenge.tex
% 20 Graduate-Level Multiple-Choice Questions on Hidden Markov Models
% Focus: Viterbi Algorithm, Baum-Welch (Forward-Backward) Algorithm
% References: CIS 5110 (Penn), Brown CS1820 lecture notes

\documentclass[12pt]{article}

\usepackage{amsmath, amssymb, amsthm}
\usepackage{geometry}
\usepackage{enumitem}
\usepackage{hyperref}
\usepackage{booktabs}
\usepackage{array}
\usepackage{xcolor}
\usepackage{parskip}

\geometry{margin=1in}

\hypersetup{
    colorlinks=true,
    linkcolor=blue,
    urlcolor=blue
}

\title{\textbf{20 Challenging Graduate-Level MCQs}\\
       \large Hidden Markov Models: Viterbi, Baum--Welch, and Beyond}
\author{HMM Study Notes\\
        \small References: CIS~5110 (UPenn), Brown CS1820}
\date{}

\begin{document}

\maketitle
\tableofcontents
\newpage

% ============================================================
\section{Introduction}
% ============================================================

This document contains \textbf{20 highly challenging graduate-level multiple-choice
questions} on \emph{Hidden Markov Models} (HMMs), with particular emphasis on:

\begin{itemize}
    \item The \textbf{Viterbi algorithm} for MAP state-sequence decoding,
    \item The \textbf{Baum--Welch algorithm} (a special case of Expectation--Maximization)
          for unsupervised parameter estimation,
    \item The \textbf{Forward--Backward algorithm} and its connection to the E-step of
          Baum--Welch,
    \item Complexity, convergence, and subtle theoretical properties.
\end{itemize}

The questions are drawn from topics covered in advanced probabilistic machine-learning
courses, notably the \textbf{CIS~5110} notes (University of Pennsylvania) and the
\textbf{Brown CS1820} lecture notes on sequence models.  They assume familiarity with
dynamic programming, EM theory, and standard probabilistic graphical model notation.

\medskip
\noindent
\textbf{Notation used throughout.}
\begin{itemize}
    \item $N$ = number of hidden states; $M$ = size of observation alphabet;
          $T$ = sequence length.
    \item $\mathbf{A} \in [0,1]^{N\times N}$: transition matrix,
          $a_{ij} = P(q_{t+1}=j \mid q_t=i)$.
    \item $\mathbf{B} \in [0,1]^{N\times M}$: emission matrix,
          $b_j(k) = P(o_t=k \mid q_t=j)$.
    \item $\boldsymbol{\pi} \in [0,1]^N$: initial-state distribution.
    \item $\lambda = (\mathbf{A}, \mathbf{B}, \boldsymbol{\pi})$: the complete HMM parameter set.
    \item $\mathbf{O} = (o_1, o_2, \dots, o_T)$: observation sequence.
    \item $\alpha_t(i) = P(o_1,\dots,o_t, q_t=i \mid \lambda)$: forward variable.
    \item $\beta_t(i) = P(o_{t+1},\dots,o_T \mid q_t=i, \lambda)$: backward variable.
    \item $\gamma_t(i) = P(q_t=i \mid \mathbf{O}, \lambda)$: state-occupancy posterior.
    \item $\xi_t(i,j) = P(q_t=i, q_{t+1}=j \mid \mathbf{O}, \lambda)$:
          transition-occupancy posterior.
    \item $\delta_t(i)$: Viterbi path probability (max-product variable).
\end{itemize}

\newpage

% ============================================================
\section{Multiple-Choice Questions}
% ============================================================

\begin{enumerate}[label=\textbf{Q\arabic*.}, leftmargin=*, itemsep=2em]

% -------------------------------------------------------
\item \textbf{Forward Variable Recurrence}

The forward variable is defined as
$\alpha_t(i) = P(o_1, \dots, o_t,\, q_t = i \mid \lambda)$.
Which of the following correctly gives the recurrence for $\alpha_{t+1}(j)$?

\begin{enumerate}[label=\Alph*.]
    \item $\displaystyle \alpha_{t+1}(j)
          = b_j(o_{t+1})\sum_{i=1}^{N} \alpha_t(i)\, a_{ij}$
    \item $\displaystyle \alpha_{t+1}(j)
          = \sum_{i=1}^{N} \alpha_t(i)\, a_{ij}\, b_i(o_{t+1})$
    \item $\displaystyle \alpha_{t+1}(j)
          = b_j(o_{t+1})\prod_{i=1}^{N} \alpha_t(i)\, a_{ij}$
    \item $\displaystyle \alpha_{t+1}(j)
          = \sum_{i=1}^{N} \alpha_t(i)\, a_{ij} + b_j(o_{t+1})$
\end{enumerate}

% -------------------------------------------------------
\item \textbf{Complexity of the Forward Algorithm}

Given an HMM with $N$ hidden states and an observation sequence of length $T$,
what is the time complexity of the forward algorithm?

\begin{enumerate}[label=\Alph*.]
    \item $O(N T)$
    \item $O(N^2 T)$
    \item $O(2^N T)$
    \item $O(N^T)$
\end{enumerate}

% -------------------------------------------------------
\item \textbf{Viterbi vs.\ Forward Recursion}

The Viterbi algorithm replaces the summation in the forward recurrence with a
maximization. Formally, let
\[
    \delta_t(i) = \max_{q_1,\dots,q_{t-1}}
        P(q_1,\dots,q_{t-1},\, q_t=i,\, o_1,\dots,o_t \mid \lambda).
\]
The correct Viterbi recurrence is:

\begin{enumerate}[label=\Alph*.]
    \item $\displaystyle \delta_{t+1}(j)
          = b_j(o_{t+1})\max_{i}\bigl[\delta_t(i)\, a_{ij}\bigr]$
    \item $\displaystyle \delta_{t+1}(j)
          = b_j(o_{t+1})\sum_{i}\bigl[\delta_t(i)\, a_{ij}\bigr]$
    \item $\displaystyle \delta_{t+1}(j)
          = \max_{i}\bigl[\delta_t(i)\, a_{ij}\bigr]$
    \item $\displaystyle \delta_{t+1}(j)
          = b_j(o_{t+1}) + \max_{i}\bigl[\delta_t(i) + \log a_{ij}\bigr]$
\end{enumerate}

% -------------------------------------------------------
\item \textbf{Viterbi Backtracking}

After computing all $\delta_t(i)$ and the backpointer array
$\psi_t(j) = \arg\max_i\bigl[\delta_{t-1}(i)\, a_{ij}\bigr]$,
the optimal state sequence $q_1^*,\dots,q_T^*$ is recovered by:

\begin{enumerate}[label=\Alph*.]
    \item Setting $q_T^* = \arg\max_i \delta_T(i)$ and then,
          for $t = T-1, \dots, 1$, $q_t^* = \psi_{t+1}(q_{t+1}^*)$.
    \item Setting $q_1^* = \arg\max_i \pi_i$ and then,
          for $t = 2, \dots, T$, $q_t^* = \arg\max_j a_{q_{t-1}^*,j}$.
    \item Setting $q_T^* = \arg\max_i \delta_T(i)$ and then,
          for $t = T, \dots, 2$, $q_{t-1}^* = \psi_T(q_T^*)$.
    \item Setting $q_t^* = \arg\max_i \gamma_t(i)$ for every $t$.
\end{enumerate}

% -------------------------------------------------------
\item \textbf{MAP vs.\ Marginal Decoding}

Let $\hat{\mathbf{q}}_{\text{MAP}} = \arg\max_{\mathbf{q}} P(\mathbf{q}\mid\mathbf{O},\lambda)$
(Viterbi output) and $\hat{q}_t^{\text{marg}} = \arg\max_i \gamma_t(i)$
(marginal/posterior decoding). Which statement is \emph{always} true?

\begin{enumerate}[label=\Alph*.]
    \item Both methods produce identical state sequences.
    \item $\hat{\mathbf{q}}_{\text{MAP}}$ maximises $P(\mathbf{O}\mid\lambda)$.
    \item The marginal-decoding sequence $(\hat{q}_1^{\text{marg}},\dots,\hat{q}_T^{\text{marg}})$
          need not correspond to any valid path with non-zero probability under the HMM.
    \item The MAP sequence always assigns higher posterior probability to each individual
          time-step than the marginal-decoding sequence.
\end{enumerate}

% -------------------------------------------------------
\item \textbf{Likelihood Computation}

Given the forward variables $\{\alpha_t(i)\}$, the total likelihood
$P(\mathbf{O}\mid\lambda)$ is computed as:

\begin{enumerate}[label=\Alph*.]
    \item $\displaystyle P(\mathbf{O}\mid\lambda)
          = \sum_{i=1}^{N} \alpha_T(i)\,\beta_T(i)$
    \item $\displaystyle P(\mathbf{O}\mid\lambda)
          = \sum_{i=1}^{N} \alpha_T(i)$
    \item $\displaystyle P(\mathbf{O}\mid\lambda)
          = \max_{i} \alpha_T(i)$
    \item $\displaystyle P(\mathbf{O}\mid\lambda)
          = \prod_{t=1}^{T}\sum_{i=1}^{N} \alpha_t(i)$
\end{enumerate}

% -------------------------------------------------------
\item \textbf{Backward Variable Definition}

The backward variable $\beta_t(i)$ is defined as
$P(o_{t+1}, o_{t+2}, \dots, o_T \mid q_t = i,\, \lambda)$.
Its recurrence (going \emph{backwards} from $T$) is:

\begin{enumerate}[label=\Alph*.]
    \item $\displaystyle \beta_t(i)
          = \sum_{j=1}^{N} a_{ij}\, b_j(o_{t+1})\,\beta_{t+1}(j)$
    \item $\displaystyle \beta_t(i)
          = \sum_{j=1}^{N} a_{ji}\, b_j(o_{t+1})\,\beta_{t+1}(j)$
    \item $\displaystyle \beta_t(i)
          = \max_{j}\bigl[a_{ij}\, b_j(o_{t+1})\,\beta_{t+1}(j)\bigr]$
    \item $\displaystyle \beta_t(i)
          = \sum_{j=1}^{N} a_{ij}\,\beta_{t+1}(j)$
\end{enumerate}

% -------------------------------------------------------
\item \textbf{State-Occupancy Posterior}

Using forward and backward variables, the posterior probability
$\gamma_t(i) = P(q_t=i\mid\mathbf{O},\lambda)$ is:

\begin{enumerate}[label=\Alph*.]
    \item $\displaystyle \gamma_t(i)
          = \frac{\alpha_t(i)\,\beta_t(i)}{\displaystyle\sum_{j=1}^{N}\alpha_t(j)\,\beta_t(j)}$
    \item $\displaystyle \gamma_t(i)
          = \alpha_t(i)\,\beta_t(i)$
    \item $\displaystyle \gamma_t(i)
          = \frac{\alpha_t(i)}{\displaystyle\sum_{j=1}^{N}\alpha_t(j)}$
    \item $\displaystyle \gamma_t(i)
          = \frac{\alpha_t(i) + \beta_t(i)}{P(\mathbf{O}\mid\lambda)}$
\end{enumerate}

% -------------------------------------------------------
\item \textbf{Transition-Occupancy Posterior}

The joint posterior $\xi_t(i,j) = P(q_t=i,\,q_{t+1}=j\mid\mathbf{O},\lambda)$ equals:

\begin{enumerate}[label=\Alph*.]
    \item $\displaystyle \xi_t(i,j)
          = \frac{\alpha_t(i)\,a_{ij}\,b_j(o_{t+1})\,\beta_{t+1}(j)}{P(\mathbf{O}\mid\lambda)}$
    \item $\displaystyle \xi_t(i,j)
          = \alpha_t(i)\,a_{ij}\,b_j(o_{t+1})\,\beta_{t+1}(j)$
    \item $\displaystyle \xi_t(i,j)
          = \frac{\alpha_t(i)\,a_{ij}\,\beta_{t+1}(j)}{P(\mathbf{O}\mid\lambda)}$
    \item $\displaystyle \xi_t(i,j)
          = \gamma_t(i)\,\gamma_{t+1}(j)$
\end{enumerate}

% -------------------------------------------------------
\item \textbf{Baum--Welch: Re-estimation of Transition Probabilities}

In the M-step of Baum--Welch, the re-estimated transition probability $\hat{a}_{ij}$ is:

\begin{enumerate}[label=\Alph*.]
    \item $\displaystyle \hat{a}_{ij}
          = \frac{\displaystyle\sum_{t=1}^{T-1}\xi_t(i,j)}
                 {\displaystyle\sum_{t=1}^{T-1}\gamma_t(i)}$
    \item $\displaystyle \hat{a}_{ij}
          = \frac{\displaystyle\sum_{t=1}^{T}\xi_t(i,j)}
                 {\displaystyle\sum_{t=1}^{T}\gamma_t(i)}$
    \item $\displaystyle \hat{a}_{ij}
          = \frac{\displaystyle\sum_{t=1}^{T-1}\gamma_t(i)}
                 {\displaystyle\sum_{t=1}^{T-1}\xi_t(i,j)}$
    \item $\displaystyle \hat{a}_{ij}
          = \frac{\displaystyle\sum_{t=1}^{T}\xi_t(i,j)}
                 {\displaystyle\sum_{t=1}^{T}\gamma_t(i)}$
\end{enumerate}

% -------------------------------------------------------
\item \textbf{Baum--Welch: Re-estimation of Emission Probabilities}

For a discrete HMM with observation alphabet $\mathcal{V} = \{v_1,\dots,v_M\}$,
the Baum--Welch re-estimate of $b_j(k) = P(o_t = v_k \mid q_t = j)$ is:

\begin{enumerate}[label=\Alph*.]
    \item $\displaystyle \hat{b}_j(k)
          = \frac{\displaystyle\sum_{t=1}^{T}\mathbf{1}[o_t=v_k]\,\gamma_t(j)}
                 {\displaystyle\sum_{t=1}^{T}\gamma_t(j)}$
    \item $\displaystyle \hat{b}_j(k)
          = \frac{\displaystyle\sum_{t=1}^{T}\mathbf{1}[o_t=v_k]\,\gamma_t(j)}
                 {\displaystyle\sum_{k'=1}^{M}\sum_{t=1}^{T}\mathbf{1}[o_t=v_{k'}]}$
    \item $\displaystyle \hat{b}_j(k)
          = \frac{\displaystyle\sum_{t=1}^{T}\gamma_t(j)}
                 {\displaystyle\sum_{t=1}^{T}\mathbf{1}[o_t=v_k]\,\gamma_t(j)}$
    \item $\displaystyle \hat{b}_j(k)
          = \mathbf{1}[o_t = v_k]\,\gamma_t(j)$
          for any $t$ with $o_t = v_k$
\end{enumerate}

% -------------------------------------------------------
\item \textbf{Convergence Guarantees of Baum--Welch}

Which statement best characterizes the convergence properties of the Baum--Welch
algorithm?

\begin{enumerate}[label=\Alph*.]
    \item It converges to the global maximum of $P(\mathbf{O}\mid\lambda)$
          regardless of initialization.
    \item Each iteration is guaranteed to not decrease $P(\mathbf{O}\mid\lambda)$,
          but convergence is only to a local maximum or saddle point.
    \item It converges to the global maximum if $N \geq 2$.
    \item It minimizes $P(\mathbf{O}\mid\lambda)$ by construction, acting as
          a gradient descent procedure.
\end{enumerate}

% -------------------------------------------------------
\item \textbf{Log-Space Computation}

To avoid numerical underflow when computing forward/backward variables for long
sequences, practitioners work in log-space. A key operation needed is computing
$\log\!\sum_i \exp(x_i)$ from $\{x_i\}$ (the log-sum-exp trick). Why is the
following \emph{naive} formula problematic?
\[
    \log\!\sum_i \exp(x_i)
\]

\begin{enumerate}[label=\Alph*.]
    \item The exponential can overflow or underflow in floating-point arithmetic
          before the sum is taken.
    \item The logarithm is undefined for negative values of $x_i$.
    \item Dynamic programming cannot be applied in log-space.
    \item Maximization cannot be converted to log-space.
\end{enumerate}

% -------------------------------------------------------
\item \textbf{Multiple Observation Sequences}

When training an HMM on a set of $K$ independent observation sequences
$\{\mathbf{O}^{(1)}, \dots, \mathbf{O}^{(K)}\}$, the correct M-step update for
$\hat{a}_{ij}$ is:

\begin{enumerate}[label=\Alph*.]
    \item $\displaystyle \hat{a}_{ij}
          = \frac{\displaystyle\sum_{k=1}^{K}\sum_{t=1}^{T_k-1}\xi_t^{(k)}(i,j)}
                 {\displaystyle\sum_{k=1}^{K}\sum_{t=1}^{T_k-1}\gamma_t^{(k)}(i)}$
    \item $\displaystyle \hat{a}_{ij}
          = \frac{1}{K}\sum_{k=1}^{K}\hat{a}_{ij}^{(k)}$
          where $\hat{a}_{ij}^{(k)}$ is the single-sequence estimate.
    \item $\displaystyle \hat{a}_{ij}
          = \frac{\displaystyle\sum_{k=1}^{K}\xi_{T_k}^{(k)}(i,j)}
                 {\displaystyle\sum_{k=1}^{K}\gamma_{T_k}^{(k)}(i)}$
    \item Both A and B produce identical results in general.
\end{enumerate}

% -------------------------------------------------------
\item \textbf{HMM as a Latent-Variable Model}

In the EM framework applied to HMMs, the complete-data log-likelihood
$\log P(\mathbf{O}, \mathbf{Q} \mid \lambda)$ decomposes as:

\begin{enumerate}[label=\Alph*.]
    \item $\displaystyle
          \log \pi_{q_1}
          + \sum_{t=1}^{T-1}\log a_{q_t q_{t+1}}
          + \sum_{t=1}^{T}\log b_{q_t}(o_t)$
    \item $\displaystyle
          \prod_{t=1}^{T} \pi_{q_t}\, a_{q_t q_{t+1}}\, b_{q_t}(o_t)$
    \item $\displaystyle
          \sum_{t=1}^{T}\log \pi_{q_t}
          + \sum_{t=1}^{T-1}\log a_{q_t q_{t+1}}
          + \sum_{t=1}^{T}\log b_{q_t}(o_t)$
    \item $\displaystyle
          \sum_{t=1}^{T}\log \pi_{q_t}
          + \sum_{t=1}^{T}\log a_{q_t q_{t+1}}
          + \sum_{t=1}^{T}\log b_{q_t}(o_t)$
\end{enumerate}

% -------------------------------------------------------
\item \textbf{Identifiability of HMMs}

Regarding parameter identifiability of HMMs, which statement is most accurate?

\begin{enumerate}[label=\Alph*.]
    \item HMM parameters $(\mathbf{A}, \mathbf{B}, \boldsymbol{\pi})$ are generically
          identifiable up to label permutation of the hidden states.
    \item HMMs are never identifiable because any permutation of states yields
          the same likelihood for all observation sequences.
    \item Identifiability holds only when the transition matrix $\mathbf{A}$ is diagonal.
    \item HMM parameters are identifiable if and only if $N < M$.
\end{enumerate}

% -------------------------------------------------------
\item \textbf{Viterbi Algorithm in Log-Space}

When the Viterbi algorithm is implemented in log-space (to prevent underflow),
the recurrence becomes:

\begin{enumerate}[label=\Alph*.]
    \item $\displaystyle
          \log\delta_{t+1}(j)
          = \log b_j(o_{t+1})
            + \max_{i}\bigl[\log\delta_t(i) + \log a_{ij}\bigr]$
    \item $\displaystyle
          \log\delta_{t+1}(j)
          = \log b_j(o_{t+1})
            + \sum_{i}\bigl[\log\delta_t(i) + \log a_{ij}\bigr]$
    \item $\displaystyle
          \log\delta_{t+1}(j)
          = \log b_j(o_{t+1})
            \cdot \max_{i}\bigl[\log\delta_t(i) + \log a_{ij}\bigr]$
    \item $\displaystyle
          \log\delta_{t+1}(j)
          = \log\Bigl(b_j(o_{t+1}) + \max_{i}\bigl[\delta_t(i)\, a_{ij}\bigr]\Bigr)$
\end{enumerate}

% -------------------------------------------------------
\item \textbf{Scaling in the Forward Algorithm}

A standard technique to prevent underflow is to scale each forward vector
$\boldsymbol{\alpha}_t$ by a factor $c_t$ such that $\sum_i \hat{\alpha}_t(i) = 1$,
where $\hat{\alpha}_t(i) = \alpha_t(i) / c_t$. The log-likelihood is then
recovered as:

\begin{enumerate}[label=\Alph*.]
    \item $\displaystyle \log P(\mathbf{O}\mid\lambda)
          = \sum_{t=1}^{T} \log c_t$
    \item $\displaystyle \log P(\mathbf{O}\mid\lambda)
          = \prod_{t=1}^{T} c_t$
    \item $\displaystyle \log P(\mathbf{O}\mid\lambda)
          = -\sum_{t=1}^{T} \log c_t$
    \item $\displaystyle \log P(\mathbf{O}\mid\lambda)
          = \log\!\left(\sum_{t=1}^{T} c_t\right)$
\end{enumerate}

% -------------------------------------------------------
\item \textbf{Posterior Decoding and Validity}

Consider a 2-state HMM in which state 2 cannot follow state 2
(i.e., $a_{22} = 0$). Suppose that at three consecutive time steps
$t, t+1, t+2$ the marginal posterior decoder outputs $\hat{q}_t^{\text{marg}} = 2$
for all three steps. This outcome:

\begin{enumerate}[label=\Alph*.]
    \item Is impossible because the posterior marginals must satisfy all transition constraints.
    \item Can occur, because marginal decoding optimizes each time step independently
          and does not enforce transition validity.
    \item Implies $P(\mathbf{O}\mid\lambda) = 0$.
    \item Is guaranteed to also be the Viterbi path in this case.
\end{enumerate}

% -------------------------------------------------------
\item \textbf{Relation Between $\xi_t$ and $\gamma_t$}

The marginal $\gamma_t(i)$ can be derived from the joint $\xi_t(i,j)$ by:

\begin{enumerate}[label=\Alph*.]
    \item $\displaystyle \gamma_t(i) = \sum_{j=1}^{N} \xi_t(i,j)$
          for $t = 1, \dots, T-1$
    \item $\displaystyle \gamma_t(i) = \sum_{j=1}^{N} \xi_{t-1}(j,i)$
          for $t = 2, \dots, T$
    \item Both A and B are valid expressions (for their respective index ranges).
    \item $\displaystyle \gamma_t(i) = \xi_t(i,i)$
\end{enumerate}

\end{enumerate}

\newpage

% ============================================================
\section{Answer Key and Explanations}
% ============================================================

\begin{enumerate}[label=\textbf{Q\arabic*.}, leftmargin=*, itemsep=1.5em]

% -------------------------------------------------------
\item \textbf{Correct answer: A}

The forward recurrence sums over all possible previous states $i$, weights
each by the transition probability $a_{ij}$, and then multiplies by the emission
$b_j(o_{t+1})$:
\[
    \alpha_{t+1}(j)
    = b_j(o_{t+1})\sum_{i=1}^{N} \alpha_t(i)\, a_{ij}.
\]
Option B incorrectly emits from the \emph{previous} state $i$ instead of the new
state $j$.  Option C replaces summation with a product (wrong).
Option D adds instead of multiplying (incorrect factorization of joint probability).

% -------------------------------------------------------
\item \textbf{Correct answer: B}

At each of the $T$ time steps the forward variable for each of the $N$ states
is updated by summing over $N$ predecessor states.  Hence the total work is
$O(N^2 T)$.  Option A ($O(NT)$) would be correct only if transitions were
independent of previous state (impossible in a proper HMM).  Options C and D
describe exponential algorithms corresponding to naive enumeration.

% -------------------------------------------------------
\item \textbf{Correct answer: A}

The Viterbi recurrence replaces $\sum$ with $\max$:
\[
    \delta_{t+1}(j)
    = b_j(o_{t+1})\max_{i}\bigl[\delta_t(i)\, a_{ij}\bigr].
\]
Option B is the forward (sum-product) algorithm, not Viterbi.
Option C omits the emission term.
Option D is the correct \emph{log-space} version (see Q17), not the probability-space version.

% -------------------------------------------------------
\item \textbf{Correct answer: A}

Standard Viterbi backtracking:
\begin{enumerate}
    \item Find the terminal state: $q_T^* = \arg\max_i \delta_T(i)$.
    \item Recurse backwards using the stored backpointers:
          $q_t^* = \psi_{t+1}(q_{t+1}^*)$ for $t = T-1, \dots, 1$.
\end{enumerate}
Option B greedily follows the most likely transition at each step---this does not
maximize the joint path probability.
Option C incorrectly applies only the last backpointer repeatedly.
Option D is marginal (posterior) decoding, not MAP path decoding.

% -------------------------------------------------------
\item \textbf{Correct answer: C}

The Viterbi output is the \emph{joint} MAP sequence; the marginal decoder
independently maximizes $\gamma_t(i)$ at each step.  Because the $\gamma_t$
marginals are computed independently of transition constraints, the sequence of
marginal argmaxes \emph{need not be a valid path}
(e.g., it may place probability mass on a transition with $a_{ij}=0$).
Option A is false: the two methods can disagree.
Option B is false: Viterbi maximises $P(\mathbf{Q}\mid\mathbf{O},\lambda)$,
not $P(\mathbf{O}\mid\lambda)$.
Option D is false: there is no such dominance guarantee per time step.

% -------------------------------------------------------
\item \textbf{Correct answer: B}

At time $T$ all future observations have been consumed, so the backward variable
$\beta_T(i) = 1$ for all $i$, giving
$\alpha_T(i)\,\beta_T(i) = \alpha_T(i)$.  Hence:
\[
    P(\mathbf{O}\mid\lambda) = \sum_{i=1}^{N} \alpha_T(i).
\]
Option A is also correct as a formula (since $\beta_T(i)=1$), but it is
unnecessarily verbose; if used with $t < T$ one must use both $\alpha$ and
$\beta$.  As stated, the simplest and most direct formula is B.

\textit{Note:} A is a valid expression at any $t$:
$\sum_i \alpha_t(i)\beta_t(i) = P(\mathbf{O}\mid\lambda)$ for all $t$.
If the question asks for the \emph{terminal} expression, B is the cleanest.

% -------------------------------------------------------
\item \textbf{Correct answer: A}

The backward recurrence proceeds from $T$ to $1$:
\[
    \beta_t(i) = \sum_{j=1}^{N} a_{ij}\, b_j(o_{t+1})\,\beta_{t+1}(j).
\]
Option B has the transition matrix transposed ($a_{ji}$ instead of $a_{ij}$).
Option C uses a max instead of sum.
Option D omits the emission term $b_j(o_{t+1})$.

% -------------------------------------------------------
\item \textbf{Correct answer: A}

By definition $\gamma_t(i) = P(q_t=i\mid\mathbf{O},\lambda)$. Bayes' theorem gives:
\[
    \gamma_t(i)
    = \frac{P(q_t=i,\mathbf{O}\mid\lambda)}{P(\mathbf{O}\mid\lambda)}
    = \frac{\alpha_t(i)\,\beta_t(i)}
           {\displaystyle\sum_{j=1}^{N}\alpha_t(j)\,\beta_t(j)}.
\]
Option B omits the normalisation.
Option C uses only the forward variable (ignores future observations).
Option D adds $\alpha$ and $\beta$, which is dimensionally and probabilistically incorrect.

% -------------------------------------------------------
\item \textbf{Correct answer: A}

\[
    \xi_t(i,j)
    = \frac{P(q_t=i,\,q_{t+1}=j,\,\mathbf{O}\mid\lambda)}{P(\mathbf{O}\mid\lambda)}
    = \frac{\alpha_t(i)\,a_{ij}\,b_j(o_{t+1})\,\beta_{t+1}(j)}
           {P(\mathbf{O}\mid\lambda)}.
\]
Option B omits the normalisation denominator.
Option C omits the emission $b_j(o_{t+1})$.
Option D would be correct only if $q_t$ and $q_{t+1}$ were independent given $\mathbf{O}$,
which they are not (due to the Markov structure).

% -------------------------------------------------------
\item \textbf{Correct answer: A}

The expected number of transitions from $i$ to $j$ is
$\sum_{t=1}^{T-1}\xi_t(i,j)$, and the expected number of times the chain leaves
state $i$ is $\sum_{t=1}^{T-1}\gamma_t(i)$:
\[
    \hat{a}_{ij}
    = \frac{\displaystyle\sum_{t=1}^{T-1}\xi_t(i,j)}
           {\displaystyle\sum_{t=1}^{T-1}\gamma_t(i)}.
\]
The upper limit is $T-1$ (not $T$) because no transition originates from the
last time step; options B and D extend both sums to $t=T$, which introduces a
systematic bias by including a time point where $\xi_t(i,j)$ is undefined.
Option C inverts numerator and denominator.

% -------------------------------------------------------
\item \textbf{Correct answer: A}

The fraction counts the expected number of times state $j$ emits symbol $v_k$
divided by the total expected time spent in state $j$:
\[
    \hat{b}_j(k)
    = \frac{\displaystyle\sum_{t:\,o_t=v_k}\gamma_t(j)}
           {\displaystyle\sum_{t=1}^{T}\gamma_t(j)}.
\]
Option B normalises by the count of symbol $v_k$ regardless of state---wrong.
Option C inverts the ratio.
Option D is not a proper formula (depends on the choice of a single $t$).

% -------------------------------------------------------
\item \textbf{Correct answer: B}

Baum--Welch is an EM algorithm applied to HMMs.  EM guarantees that the
observed-data log-likelihood $\log P(\mathbf{O}\mid\lambda)$ is non-decreasing
at each iteration.  However, the objective is non-convex in $\lambda$, so the
algorithm can converge to a \emph{local} maximum or saddle point depending on
initialisation.  Global convergence (option A) is not guaranteed.
Option C has no theoretical basis.
Option D is incorrect; EM monotonically \emph{increases} (not decreases) the likelihood.

% -------------------------------------------------------
\item \textbf{Correct answer: A}

Computing $\exp(x_i)$ for very negative $x_i$ (e.g.\ $x_i = -1000$) yields
machine zero (underflow), and for very positive values it overflows.
The standard log-sum-exp trick subtracts $\max_i x_i$ before exponentiating:
\[
    \log\!\sum_i e^{x_i}
    = m + \log\!\sum_i e^{x_i - m}, \quad m = \max_i x_i.
\]
Option B is irrelevant: $x_i = \log\alpha_t(i) \leq 0$ in practice, but the
$\log$ is of a \emph{sum of exponentials}, not of $x_i$ directly.
Options C and D are false claims.

% -------------------------------------------------------
\item \textbf{Correct answer: A}

With $K$ independent sequences the complete-data sufficient statistics are
summed across sequences, and the M-step re-estimate becomes:
\[
    \hat{a}_{ij}
    = \frac{\displaystyle\sum_{k=1}^{K}\sum_{t=1}^{T_k-1}\xi_t^{(k)}(i,j)}
           {\displaystyle\sum_{k=1}^{K}\sum_{t=1}^{T_k-1}\gamma_t^{(k)}(i)}.
\]
Option B (simple average of per-sequence estimates) is incorrect unless all
sequences have equal length and equal posterior mass on state $i$;
options A and B are generally different.
Option C sums only the \emph{final} time step of each sequence, which is incorrect.

% -------------------------------------------------------
\item \textbf{Correct answer: A}

The joint distribution factorises as:
\[
    P(\mathbf{O},\mathbf{Q}\mid\lambda)
    = \pi_{q_1}\prod_{t=1}^{T-1}a_{q_t q_{t+1}}\prod_{t=1}^{T}b_{q_t}(o_t).
\]
Taking $\log$:
\[
    \log P(\mathbf{O},\mathbf{Q}\mid\lambda)
    = \log\pi_{q_1}
      + \sum_{t=1}^{T-1}\log a_{q_t q_{t+1}}
      + \sum_{t=1}^{T}\log b_{q_t}(o_t).
\]
Option A is correct.
Option B is in probability space (not log space).
Option C incorrectly includes $\log\pi_{q_t}$ for every $t$ rather than only $t=1$.
Option D incorrectly extends the transition sum to $t=T$ (there is no state $q_{T+1}$) and also sums $\log\pi$ at every step.

% -------------------------------------------------------
\item \textbf{Correct answer: A}

A seminal result (Allman, Matias \& Rhodes, 2009) shows that HMM parameters are
\emph{generically identifiable} up to permutation of the hidden states,
provided the emission distributions are linearly independent.
Option B would be correct only for degenerate (permutation-equivalent) models,
and identifiability up to permutation is standard.
Option C is false.
Option D ($N < M$) is a \emph{sufficient} condition in some settings but not
necessary in general.

% -------------------------------------------------------
\item \textbf{Correct answer: A}

Taking $\log$ of the Viterbi recurrence $\delta_{t+1}(j) = b_j(o_{t+1})\max_i[\delta_t(i)a_{ij}]$ and using $\log(xy) = \log x + \log y$:
\[
    \log\delta_{t+1}(j)
    = \log b_j(o_{t+1}) + \max_{i}\bigl[\log\delta_t(i) + \log a_{ij}\bigr].
\]
Option B sums instead of taking max---this is the log-space forward algorithm.
Option C multiplies instead of adding logs (incorrect).
Option D takes the log after mixing probability-space and log-space terms (incorrect).

% -------------------------------------------------------
\item \textbf{Correct answer: A}

Define $c_t = \sum_i \alpha_t(i)$ (the normalisation constant at step $t$).
The scaled forward variable is $\hat{\alpha}_t(i) = \alpha_t(i)/\prod_{s=1}^t c_s$.
The total likelihood equals the product of all scaling constants:
\[
    P(\mathbf{O}\mid\lambda) = \prod_{t=1}^{T} c_t,
\]
so $\log P(\mathbf{O}\mid\lambda) = \sum_{t=1}^{T}\log c_t$.
Option B returns the unlogged product (not log-likelihood).
Option C negates the sum (would be $-\log P$).
Option D takes the log of the sum of $c_t$, which is incorrect.

% -------------------------------------------------------
\item \textbf{Correct answer: B}

Marginal decoding maximises $\gamma_t(i)$ independently at each time step
without enforcing transition constraints.  If $a_{22}=0$, the Viterbi
algorithm would never place state 2 at two consecutive time steps, but the
posterior marginals $\gamma_t(2)$ can still be the largest marginals at
multiple consecutive steps (because they integrate over \emph{all} paths, not
just those through state~2).  Thus option B is correct.
Option A is false.
Option C is false: a positive observation likelihood is possible even if the
most probable \emph{marginal} state sequence is invalid.
Option D is false: the Viterbi path would never repeat state~2 here.

% -------------------------------------------------------
\item \textbf{Correct answer: C}

Both directions are valid for the index ranges given:
\begin{align*}
    \gamma_t(i) &= \sum_{j=1}^{N} \xi_t(i,j),
    \quad t = 1, \dots, T-1 \quad \text{(marginalise over next state)}, \\
    \gamma_t(i) &= \sum_{j=1}^{N} \xi_{t-1}(j,i),
    \quad t = 2, \dots, T \quad \text{(marginalise over previous state)}.
\end{align*}
Therefore option C (``both A and B'') is correct.
Option D, $\gamma_t(i) = \xi_t(i,i)$, equates the marginal to the
self-transition occupancy, which is wrong unless $N=1$.

\end{enumerate}

\newpage

% ============================================================
\section*{Quick Reference: Key Formulas}
% ============================================================
\addcontentsline{toc}{section}{Quick Reference: Key Formulas}

\begin{center}
\renewcommand{\arraystretch}{2.0}
\begin{tabular}{lp{9cm}}
\toprule
\textbf{Quantity} & \textbf{Formula} \\
\midrule
Forward init & $\alpha_1(i) = \pi_i\, b_i(o_1)$ \\
Forward recur & $\alpha_{t+1}(j) = b_j(o_{t+1})\sum_i \alpha_t(i)\,a_{ij}$ \\
Backward init & $\beta_T(i) = 1$ \\
Backward recur & $\beta_t(i) = \sum_j a_{ij}\,b_j(o_{t+1})\,\beta_{t+1}(j)$ \\
Likelihood & $P(\mathbf{O}\mid\lambda) = \sum_i \alpha_T(i)$ \\
State posterior & $\gamma_t(i) = \dfrac{\alpha_t(i)\beta_t(i)}{\sum_j \alpha_t(j)\beta_t(j)}$ \\
Trans.\ posterior & $\xi_t(i,j) = \dfrac{\alpha_t(i)\,a_{ij}\,b_j(o_{t+1})\,\beta_{t+1}(j)}{P(\mathbf{O}\mid\lambda)}$ \\
BW: $\hat{a}_{ij}$ & $\dfrac{\sum_{t=1}^{T-1}\xi_t(i,j)}{\sum_{t=1}^{T-1}\gamma_t(i)}$ \\
BW: $\hat{b}_j(k)$ & $\dfrac{\sum_t \mathbf{1}[o_t=v_k]\,\gamma_t(j)}{\sum_t \gamma_t(j)}$ \\
Viterbi recur & $\delta_{t+1}(j) = b_j(o_{t+1})\max_i[\delta_t(i)\,a_{ij}]$ \\
\bottomrule
\end{tabular}
\end{center}

\end{document}
